\section{Diagonal random-phase form factor at times one, two and three}
\label{sec:diagonal-phase-form-factor}
The phases below are independent at the \(N\) sites of the position basis.
They are not the sixteen mode-pair phases of the kicked ensemble, and the
propagator is not the fixed kick. The canonical proof is
\path{mathematics/fixed-map/diagonal_phase_form_factor.md}.
The integer comparisons are certified by
\path{formalization/Bgs/DiagonalPhaseFormFactorRational.lean} and recomputed
by the Python and C++ verifiers named there. No external refereeing is asserted.
BGS-TRACE-005 and BGS-BGS-001 remain open.

\begin{theorem}[Diagonal random-phase form factor, BGS-TRACE-127]
\label{thm:diagonal-phase-form-factor}
Let \(M\) be an \(N\times N\) complex matrix with every entry of modulus
\(N^{-1/2}\), and let \(D\) be diagonal with independent uniform phases.
For \(U=DM\) and every integer \(m\ge 1\),
\[
\mathbb E\lvert\operatorname{Tr} U^m\rvert^2\le m!.
\]
The first two times are exact:
\(\mathbb E\lvert\operatorname{Tr} U\rvert^2=1\) and
\(\mathbb E\lvert\operatorname{Tr} U^2\rvert^2=(2N-1)/N\).
For the cat-map quantization at every odd \(N\ge 3\),
\[
\mathbb E\lvert\operatorname{Tr} U_N^3\rvert^2=\frac{6N^2-9N+4}{N^2},
\]
which lies strictly between \(1\) and \(6\). On the same dimensions,
\(\lvert\ker(S^2-I)\rvert=1\) and \(\lvert\ker(S^3-I)\rvert\) equals \(49\)
if \(7\) divides \(N\) and equals \(1\) otherwise. By the trace-closure lemma
of BGS-TRACE-126 these are the pure squared traces, so the diagonal-phase
value lies above the pure value at time \(2\), above it at time \(3\) when
\(7\) does not divide \(N\), and below \(6\) while the pure value equals
\(49\) when \(7\) divides \(N\).

The time-three limit \(6\) is not the CUE value \(3\). No symmetry class is
identified. The bound \(m!\) does not separate a revival once \(m!\ge N^2\).
Nothing beyond that bound is claimed for \(m\ge 4\). A deterministic
twenty-four-draw sample of the sixteen-phase ensemble, recomputed in both
verifiers, has mean \(44.632431403503\) at \(N=7\), time \(3\), \(\eta=2\),
and mean \(47.417022055369\) at \(N=31\), time \(5\), \(\eta=2\), where the
kernel of \(S^5-I\) has order \(1\). That sample is not an evaluation of
the sixteen-phase expectation and not a CUE or COE limit.
\end{theorem}
